%% Version: 0.8 (12.02.2019)

%% My_document_info.tex
%% Copyright 2020 KIT, IAR-IPR (Denis Štogl)
%% Mail: denis.stogl@kit.edu
%
% This work may be distributed and/or modified under the
% conditions of the LaTeX Project Public License, either version 1.3
% of this license or (at your option) any later version.
% The latest version of this license is in
%   http://www.latex-project.org/lppl.txt
% and version 1.3 or later is part of all distributions of LaTeX
% version 2005/12/01 or later.
%
% This work has the LPPL maintenance status `maintained'.
%
% The Current Maintainer of this work is M. Y. Name.
%
% This work consists of the files pig.dtx and pig.ins
% and the derived file thesis.tex.


\title{Vergleich von Distributed-Ledger-Architekturen für das dezentrale Datenmanagement im Gesundheitswesen}
\titleotherlanguage{Comparison of distributed ledger architectures for decentralized data management in healthcare}

\author{Arsenii Tochenov}
\address{Kaiserstr. 191}
\city{76133 Karlsruhe}
\email{ubvaa@student.kit.edu}

\keywords{Self-Sovereign Identity, Verifiable Data Registry, Distributed Ledger Technology, Hashgraph, Hedera Consensus Service, Hyperledger Indy, AnonCreds, Decentralized Identifiers, Elektronische Gesundheitsakte, Performanzevaluation, Skalierbarkeit}
\keywordsotherlanguge{Self-Sovereign Identity, Verifiable Data Registry, Distributed Ledger Technology, Hashgraph, Hedera Consensus Service, Hyperledger Indy, AnonCreds, Decentralized Identifiers, Electronic Health Records, Performance Evaluation, Scalability}

%% Study program or a seminar/subject
%% \studyprogram{Intelligente Industrieroboter}

%% IMPORTANT: Seminar only: If not "Seminar Intelligente Industrieroboter" uncomment this
\nogrouplogo
\noinstitutelogo

%% Name of your institute (Default: IAR-IPR)
\institute{ITIV}
%% Name of your faculty (Default: KIT-Fakultät für Informatik
\KITfaculty{ Institut für Technik der Informationsverarbeitung}
%% Address of your institute (Default: Engler-Bunte-Ring 8)
\instituteaddress{Engesserstraße 5 (Geb. 30.10)}
% %% Insitute City (Default: 76131 Karlsruhe)
% \institutecity{}

\reviewerone{Prof. Dr. Ralf Reussner (KASTEL)}
\reviewertwo{Prof. Dr. Wilhelm Stork (ITIV)}
%
% %% The advisors are PhDs or Postdocs
\advisorone{Dr.-Ing. Christina Erler}
% %% The second advisor can be omitted
\advisortwo{M.Sc. Florian Mazura}
%
% %% Please enter the start end end time of your thesis (for techreport not needed)
\editingtime{26. Mai 2026}{28. September 2026}

%% --------------------------------
%% | Settings for word separation |
%% --------------------------------
% Help for separation:
% In german package the following hints are additionally available:
% "- = Additional separation
% "| = Suppress ligation and possible separation (e.g. Schaf"|fell)
% "~ = Hyphenation without separation (e.g. bergauf und "~ab)
% "= = Hyphenation with separation before and after
% "" = Separation without a hyphenation (e.g. und/""oder)

% Describe separation hints here:
\hyphenation{
% Pro-to-koll-in-stan-zen
% Ma-na-ge-ment  Netz-werk-ele-men-ten
% Netz-werk Netz-werk-re-ser-vie-rung
% Netz-werk-adap-ter Fein-ju-stier-ung
% Da-ten-strom-spe-zi-fi-ka-tion Pa-ket-rumpf
% Kon-troll-in-stanz
}

%%
%% --------------------
%% |   Bibliography   |
%% --------------------
\newcommand{\mybibliographyfiles}{Bibliography/my_thesis_bibliography}


%% --------------------
%% |     Acronyms     |
%% --------------------
%\newacronym{ipr}{IAR-IPR}{Institute for Anthropomatics and Robotics - Intelligent %Process Control and Robotics}

%% --------------------
%% |     Glossary     |
%% --------------------
%\newglossaryentry{robot}
%{
%    name=robot,
%    description={The robot developed in this work.}
%}